\documentclass[11pt]{article}

% ============================================================
% PACKAGES
% ============================================================

\usepackage[margin=1in]{geometry}
\usepackage{amsmath,amssymb}
\usepackage{array}
\usepackage{booktabs}
\usepackage{tabularx}
\usepackage{enumitem}
\usepackage{microtype}
\usepackage[most]{tcolorbox}
\usepackage{tikz}
\usepackage{hyperref}

% ============================================================
% FORMATTING
% ============================================================

\setlength{\parindent}{0pt}
\setlength{\parskip}{0.65em}

\setlist[itemize]{itemsep=0.3em, topsep=0.4em}
\setlist[enumerate]{itemsep=0.3em, topsep=0.4em}

\hypersetup{
    colorlinks=false,
    pdfborder={0 0 0}
}

% ============================================================
% WHITEBOARD ENVIRONMENT
% ============================================================

\newtcolorbox{whiteboard}{
    enhanced,
    breakable,
    colback=white,
    colframe=black,
    boxrule=0.9pt,
    sharp corners,
    title=\textbf{WHITEBOARD},
    fonttitle=\bfseries,
    attach boxed title to top left={
        xshift=0.4cm,
        yshift=-2mm
    },
    boxed title style={
        colback=white,
        colframe=black,
        sharp corners
    },
    before skip=1em,
    after skip=1em
}

\newcommand{\boardtakeaway}[1]{
    \begin{center}
    \fbox{
        \parbox{0.88\linewidth}{
            \centering
            \textbf{#1}
        }
    }
    \end{center}
}

% Modern teaching notation for exclusive disjunction.
% This is not Stoic notation.
\newcommand{\xor}{\mathbin{\mathrm{XOR}}}

% ============================================================
% DOCUMENT
% ============================================================

\begin{document}

\begin{center}
    {\LARGE \textbf{The Megarians and the Stoics}}\\[0.4em]
    {\large From Dialectical Paradox to Propositional Logic}\\[0.8em]
    \textit{History of Logic}\\[0.5em]
    \textbf{Primary source:} William and Martha Kneale,
    \textit{The Development of Logic}, Chapter III,
    ``The Megarians and the Stoics.''
\end{center}

\vspace{1em}

% ============================================================
% INTRODUCTION
% ============================================================

\section{Two Great Traditions of Ancient Logic}

The preceding chapter concentrated upon Aristotle.

Aristotle's mature formal theory grew chiefly out of reflection
upon reasoning involving general terms. Geometry and the ideal of
demonstration provided an especially important model.

But Aristotle's was not the only great logical tradition of the
ancient world.

Later antiquity distinguished two major schools:

\begin{enumerate}
    \item the \textbf{Peripatetic} tradition deriving from Aristotle;
    \item the \textbf{Stoic} tradition, developed especially by
    Chrysippus from Megarian logical work.
\end{enumerate}

The difference between the two traditions is not simply a
difference of doctrine.

They begin from different kinds of argumentative practice and
therefore isolate different kinds of logical structure.

\begin{whiteboard}
\textbf{TWO GREAT ANCIENT TRADITIONS}

\[
\begin{array}{ccc}
\textbf{PERIPATETIC}
&&
\textbf{MEGARIAN--STOIC}
\\[0.4em]
\text{Aristotle}
&&
\text{Megarians}
\\
\downarrow
&&
\downarrow
\\
\text{term / syllogistic logic}
&&
\text{Chrysippus}
\\
&&
\downarrow
\\
&&
\text{propositional / dialectical logic}
\end{array}
\]

\boardtakeaway{Independent traditions, but ultimately complementary.}
\end{whiteboard}

Unfortunately, the writings of the great Stoic logicians have
not survived intact.

Our knowledge comes primarily from later reports, often written
by people who belonged to rival philosophical schools.

We can reconstruct the major features of Stoic logic, but the
evidence is sometimes incomplete or uncertain.

% ============================================================
\section{Geometry and Dialectic: Two Sources of Logic}
% ============================================================

The contrast between Aristotle and the Megarians can be traced
back to the kinds of reasoning upon which they concentrated.

Aristotle appears to have been strongly influenced by
demonstrative reasoning of the sort found in mathematics.

The Megarians, by contrast, inherited traditions associated with:

\begin{itemize}
    \item Eleatic philosophy;
    \item Zeno's dialectical arguments;
    \item reductio;
    \item eristic disputation;
    \item paradox;
    \item conditional argument.
\end{itemize}

These different origins help explain why Aristotle and the Stoics
developed such different formal systems.

\begin{whiteboard}
\textbf{TWO SOURCES $\rightarrow$ TWO LOGICS}

\[
\begin{array}{ccc}
\textbf{GEOMETRY}
&&
\textbf{DIALECTIC}
\\
\downarrow
&&
\downarrow
\\
\text{Aristotle}
&&
\text{Megarians / Stoics}
\\
\downarrow
&&
\downarrow
\\
\text{general terms}
&&
\text{whole propositions}
\\
\text{syllogisms}
&&
\text{connectives / inference}
\end{array}
\]

\boardtakeaway{Different argumentative practices reveal different logical structures.}
\end{whiteboard}

This contrast will organize the entire lecture.

% ============================================================
% PART I
% ============================================================

\part*{I. The Megarian Tradition and Its Paradoxes}

% ============================================================
\section{From Euclides to Chrysippus}
% ============================================================

Euclides of Megara, an older contemporary of Plato, is
traditionally regarded as the founder of the Megarian school.

The line of influence is not perfectly recoverable, but several
figures are especially important.

Eubulides became famous for paradoxes.

Diodorus Cronus and Philo developed theories of modality and
conditionals.

Stilpo belonged to the Megarian tradition and taught Zeno of
Citium, the founder of Stoicism.

Chrysippus, the third major head of the Stoic school, eventually
gave Stoic dialectic its most systematic form.

\begin{whiteboard}
\textbf{MEGARIAN $\rightarrow$ STOIC LINE}

\[
\begin{array}{c}
\text{Euclides of Megara}\\
\downarrow\\
\text{Eubulides / Megarian tradition}\\
\downarrow\\
\text{Diodorus Cronus + Philo}\\
\downarrow\\
\text{Zeno of Citium}\\
\downarrow\\
\boxed{\text{Chrysippus}}
\end{array}
\]

\[
\text{Chrysippus}
=
\text{great systematizer of Stoic logic}
\]
\end{whiteboard}

Ancient admiration for Chrysippus as a dialectician was
extraordinary.

His logical work was regarded as a rival to Aristotle's system,
not merely as a minor supplement to it.

% ============================================================
\section{Eubulides and the Logical Use of Paradox}
% ============================================================

Eubulides was traditionally credited with several paradoxes.

The Kneales reduce the traditional list to four principal types:

\begin{enumerate}
    \item the Liar;
    \item the Hooded Man, also called the Unnoticed Man or Electra;
    \item the Bald Man or Heap;
    \item the Horned Man.
\end{enumerate}

These puzzles are important because each exposes a genuine
logical or philosophical problem.

They should not be dismissed as mere verbal games.

\begin{whiteboard}
\textbf{EUBULIDES: FOUR PARADOXES}

\[
\begin{array}{ll}
1. & \text{Liar}\\
2. & \text{Hooded Man / Electra}\\
3. & \text{Bald Man / Heap}\\
4. & \text{Horned Man}
\end{array}
\]

\boardtakeaway{Paradox becomes a tool for discovering logical problems.}
\end{whiteboard}

% ============================================================
\section{The Liar}
% ============================================================

The Liar concerns a statement that turns upon its own truth or
falsity.

A modernized compact form is:

\begin{center}
``This statement is false.''
\end{center}

Suppose the statement is true.

Then what it says is the case, and therefore it is false.

But suppose it is false.

Then what it says is not the case, which seems to make it true.

\begin{whiteboard}
\textbf{THE LIAR}

\[
\text{``This statement is false.''}
\]

If TRUE:

\[
\text{what it says} \Rightarrow \text{FALSE}
\]

If FALSE:

\[
\text{what it says appears} \Rightarrow \text{TRUE}
\]

\[
\boxed{\text{Problem: SELF-REFERENCE + TRUTH}}
\]
\end{whiteboard}

The purpose of this lecture is not to solve the Liar.

Its historical importance lies in the fact that Megarian
dialecticians had isolated a profound difficulty involving
self-reference and the predicates ``true'' and ``false.''

% ============================================================
\section{The Hooded Man}
% ============================================================

The Hooded Man raises a different problem.

Suppose you know your brother.

A hooded person enters.

Unknown to you, the hooded person is your brother.

You do not recognize the hooded person.

Must we therefore say both:

\begin{center}
you know your brother
\end{center}

and:

\begin{center}
you do not know your brother?
\end{center}

The puzzle concerns the interaction of identity with contexts
involving knowledge.

\begin{whiteboard}
\textbf{HOODED MAN}

Suppose:

\[
X=Y
\]

Question:

\[
\text{Can everything said truly of }X
\text{ be substituted for }Y?
\]

\vspace{0.4em}

\boardtakeaway{Identity alone may not preserve truth in knowledge-contexts.}
\end{whiteboard}

The Kneales also note that the puzzle forces us to distinguish
different uses of the word ``know.''

We need not yet develop a full theory of such contexts.

The important point is that simple substitution of identical
objects can become problematic when the proposition also
describes someone's knowledge.

% ============================================================
\section{The Heap and the Bald Man}
% ============================================================

The Heap and the Bald Man are versions of what later came to be
called the sorites paradox.

Suppose a man with only one hair is bald.

Would two hairs make him non-bald?

Three?

Four?

At what precise point does the change occur?

The heap version asks the corresponding question about grains
of material.

\begin{whiteboard}
\textbf{THE HEAP / BALD MAN}

\[
1\text{ hair} \rightarrow \text{bald?}
\]

\[
2\text{ hairs} \rightarrow \text{bald?}
\]

\[
3\text{ hairs} \rightarrow \text{bald?}
\]

\[
\vdots
\]

\[
\boxed{\text{WHERE IS THE BOUNDARY?}}
\]

\[
\text{Problem: VAGUE PREDICATES}
\]
\end{whiteboard}

The paradox reveals that many ordinary predicates do not possess
sharp boundaries of application.

% ============================================================
\section{The Horned Man}
% ============================================================

The Horned Man takes the form:

\begin{center}
What you have not lost, you still have.
\end{center}

\begin{center}
You have not lost horns.
\end{center}

Therefore:

\begin{center}
You still have horns.
\end{center}

The argument is fallacious, but the interesting question is why.

The Kneales interpret the difficulty in terms of
\textbf{presupposition}.

The statement:

\begin{center}
``You have lost your horns''
\end{center}

normally presupposes that you once had horns.

But a denial can operate in more than one way.

It may reject only the explicit claim that you lost them while
leaving the presupposition intact, or it may reject the entire
complex claim.

\begin{whiteboard}
\textbf{HORNED MAN}

\[
\text{What you have not lost, you have.}
\]

\[
\text{You have not lost horns.}
\]

\[
\therefore
\text{You have horns.}
\]

\vspace{0.5em}

Hidden issue:

\[
\text{``You lost horns''}
\]

presupposes:

\[
\text{``You once had horns.''}
\]

\boardtakeaway{Negation may reject more than the explicit predicate.}
\end{whiteboard}

% ============================================================
\section{What the Four Paradoxes Reveal}
% ============================================================

The puzzles attributed to Eubulides anticipate four enormous
areas of later logical inquiry.

\begin{whiteboard}
\textbf{PARADOX $\rightarrow$ LOGICAL PROBLEM}

\[
\begin{array}{rcl}
\text{Liar}
&\rightarrow&
\text{truth / self-reference}
\\[0.4em]
\text{Hooded Man}
&\rightarrow&
\text{identity / knowledge}
\\[0.4em]
\text{Heap}
&\rightarrow&
\text{vagueness}
\\[0.4em]
\text{Horned Man}
&\rightarrow&
\text{presupposition / negation}
\end{array}
\]
\end{whiteboard}

Whatever role entertainment and controversy played in Megarian
practice, the intellectual content of these puzzles is serious.

% ============================================================
\section{A Historical Mistake: Rivalry Instead of Combination}
% ============================================================

The Megarian and Stoic schools became involved in a long rivalry
with the Peripatetic tradition.

The Kneales regard this as unfortunate.

Aristotelian and Stoic logic are not naturally alternatives.

They study different dimensions of logical structure.

\begin{whiteboard}
\textbf{ANCIENT POLEMIC}

\[
\text{Aristotle}
\quad\mathbf{VS.}\quad
\text{Stoics}
\]

\[
\text{Better picture:}
\]

\[
\boxed{
\text{Aristotle}
+
\text{Stoics}
}
\]

\[
\text{TERM LOGIC}
+
\text{PROPOSITIONAL LOGIC}
\]
\end{whiteboard}

The failure to integrate these two achievements will become one
of the chapter's central historical conclusions.

% ============================================================
% PART II
% ============================================================

\part*{II. Megarian and Stoic Theories of Modality}

% ============================================================
\section{The Problem of Possibility}
% ============================================================

Aristotle reports that some Megarians denied a genuine
distinction between actuality and potentiality.

Such a doctrine fits naturally with an Eleatic suspicion of
motion and change.

Yet later Megarians devoted substantial attention to modal
notions such as:

\begin{itemize}
    \item possible;
    \item impossible;
    \item necessary;
    \item non-necessary.
\end{itemize}

The question becomes:

\begin{center}
Can modality be explained without Aristotle's theory of real
potentiality?
\end{center}

\begin{whiteboard}
\textbf{THE MODAL QUESTION}

What do we mean by:

\[
\begin{array}{c}
\text{POSSIBLE}\\
\text{IMPOSSIBLE}\\
\text{NECESSARY}\\
\text{NON-NECESSARY}
\end{array}
\]

\[
\boxed{\text{Diodorus} \neq \text{Philo} \neq \text{later Stoics}}
\]
\end{whiteboard}

% ============================================================
\section{Diodorus Cronus: Modal Definitions}
% ============================================================

Diodorus offers the most striking of the Megarian modal
theories.

He defines modality through the truth or falsity of a statement
across time.

Using \(P\) as a neutral placeholder for the relevant statement:

\begin{whiteboard}
\textbf{DIODORUS}

\[
\boxed{
\text{Possible }P
=
P\text{ is true or will be true}
}
\]

\[
\boxed{
\text{Impossible }P
=
P\text{ is false and will never be true}
}
\]

\[
\boxed{
\text{Necessary }P
=
P\text{ is true and will never be false}
}
\]

\[
\boxed{
\text{Non-necessary }P
=
P\text{ is false or will be false}
}
\]
\end{whiteboard}

The modal predicates are thus tied directly to patterns of truth
through time.

% ============================================================
\section{Diodoran Modality Is Modality-at-a-Time}
% ============================================================

The Kneales emphasize that Diodorus is not defining necessity
simply and timelessly.

He is effectively defining \textbf{necessity-at-a-time}.

Suppose a statement is now true and will never again become
false.

It is now necessary according to Diodorus.

But before the relevant event occurred, the corresponding
statement may merely have been possible.

Thus modality can change.

\begin{whiteboard}
\textbf{DIODORAN MODALITY}

\[
\text{modality depends upon truth through TIME}
\]

Possible changes:

\[
\text{possible}
\rightarrow
\text{necessary}
\]

\[
\text{possible}
\rightarrow
\text{impossible}
\]

But once reached:

\[
\text{necessary}
\rightarrow
\text{stays necessary}
\]

\[
\text{impossible}
\rightarrow
\text{stays impossible}
\]
\end{whiteboard}

The theory is closely connected with Diodorus' famous
\textbf{Master Argument}.

% ============================================================
\section{The Master Argument}
% ============================================================

The Master Argument is known chiefly through later reports.

It begins with three propositions which Diodorus claims are
mutually incompatible.

\begin{enumerate}
    \item Everything that is past and true is necessary.
    \item The impossible does not follow from the possible.
    \item Something which neither is nor will be true may
    nevertheless be possible.
\end{enumerate}

Diodorus accepts the first two and rejects the third.

Hence he identifies the possible with what either is or will be
true.

\begin{whiteboard}
\textbf{THE MASTER ARGUMENT}

Three claims cannot all stand:

\[
\begin{array}{ll}
1. & \text{Past truth } \rightarrow \text{ necessary}\\[0.4em]
2. & \text{Impossible cannot follow from possible}\\[0.4em]
3. & \text{Something may be possible}\\
   & \text{although it never occurs}
\end{array}
\]

Diodorus:

\[
\text{accept }1+2
\]

\[
\text{reject }3
\]

Therefore:

\[
\boxed{
\text{Possible}
=
\text{what is or will be true}
}
\]
\end{whiteboard}

This conclusion was naturally associated in antiquity with
determinism or fatalism.

% ============================================================
\section{Why the Master Argument Seems Plausible}
% ============================================================

The first premise gains much of its force from a familiar thought:

\begin{center}
What has happened cannot now be made not to have happened.
\end{center}

Even Aristotle remarks that deliberation concerns the future,
not the already completed past.

But a distinction is required.

The past may now be \textbf{unalterable} without thereby having
been \textbf{absolutely necessary} before it occurred.

\begin{whiteboard}
\textbf{WHY PREMISE 1 TEMPTS US}

\[
\text{What has happened}
\]

\[
\downarrow
\]

\[
\text{cannot now be changed}
\]

Temptation:

\[
\text{cannot now be changed}
\Rightarrow
\text{was absolutely necessary}
\]

But:

\[
\boxed{
\text{UNALTERABILITY}
\neq
\text{ABSOLUTE NECESSITY}
}
\]
\end{whiteboard}

This is closely related to the distinction between relative and
absolute necessity encountered in Aristotle.

% ============================================================
\section{The Ambiguity in the Master Argument}
% ============================================================

The Kneales identify a crucial ambiguity.

We must distinguish:

\begin{enumerate}
    \item a statement \textbf{about the past};
    \item a statement grammatically \textbf{in the past tense}.
\end{enumerate}

These are not equivalent.

A sentence can be in the past tense while its interpretation
depends upon the present context of utterance.

Likewise, a sentence with a different tense can make a claim
about an event already past.

\begin{whiteboard}
\textbf{CRUCIAL DISTINCTION}

\[
\boxed{
\text{statement ABOUT the past}
\neq
\text{statement IN the past tense}
}
\]

Underlying problem:

\[
\boxed{
\text{sentence}
\neq
\text{what the sentence expresses}
}
\]
\end{whiteboard}

The Kneales argue that the Master Argument trades upon this
ambiguity.

The problem resembles one already encountered in Aristotle's
discussion of future contingents.

Both theories suffer from insufficiently sharp separation of
sentences from the propositions or contents expressed by them.

% ============================================================
\section{Philo's Alternative Theory of Possibility}
% ============================================================

Philo rejects Diodorus' restriction of possibility to what does or
will actually occur.

For Philo, something can be possible in virtue of the intrinsic
character of the subject even when external conditions prevent
its realization.

The ancient example concerns wood at the bottom of the sea.

The wood remains the sort of thing that can burn even though
its present surroundings prevent it from burning.

\begin{whiteboard}
\textbf{PHILO}

\[
\boxed{
\text{Possible}
=
\text{admits of truth by its intrinsic nature}
}
\]

Example:

\[
\text{wood under the sea}
\]

\[
\text{intrinsically burnable}
\]

but:

\[
\text{external circumstances prevent burning}
\]

\boardtakeaway{Possible $\neq$ what actually happens.}
\end{whiteboard}

The Kneales interpret Philo as making possibility the basic modal
notion and associating it approximately with intrinsic
consistency or admissibility.

% ============================================================
\section{The Later Stoic Variant}
% ============================================================

Later reports attribute to the Stoics a more complicated theory.

The evidence is compressed and does not permit complete
certainty.

But the important distinction appears to be between:

\begin{enumerate}
    \item what something intrinsically admits;
    \item what external circumstances allow or prevent.
\end{enumerate}

Thus modal status can involve both intrinsic and circumstantial
features.

\begin{whiteboard}
\textbf{STOIC MODALITY}

\[
\boxed{
\text{INTRINSIC ADMISSIBILITY}
+
\text{EXTERNAL CIRCUMSTANCES}
}
\]

Distinguish:

\[
\text{absolute modality}
\]

from

\[
\text{relative / circumstance-dependent modality}
\]
\end{whiteboard}

Because our ancient evidence is fragmentary, this reconstruction
must remain cautious.

% ============================================================
\section{Three Ancient Strategies for Modality}
% ============================================================

The differences can now be summarized.

\begin{whiteboard}
\textbf{THREE APPROACHES}

\[
\begin{array}{rcl}
\textbf{Diodorus}
&\rightarrow&
\text{truth through time}
\\[0.5em]
\textbf{Philo}
&\rightarrow&
\text{intrinsic admissibility}
\\[0.5em]
\textbf{Stoics}
&\rightarrow&
\text{intrinsic admissibility}\\
&&+\text{ external circumstances}
\end{array}
\]

\boardtakeaway{Same modal vocabulary; different underlying theories.}
\end{whiteboard}

The modal debate prepares directly for another Megarian
controversy: the nature of conditional statements.

% ============================================================
% PART III
% ============================================================

\part*{III. The Debate on the Nature of Conditionals}

% ============================================================
\section{Why Dialecticians Care About ``If''}
% ============================================================

Conditionals were not central to Aristotle's classification of
propositions.

For a thinker working in the tradition of Zeno, however,
conditional reasoning is unavoidable.

A typical refutation has the following shape:

\[
P\rightarrow Q,
\]

\[
P\rightarrow \neg Q,
\]

therefore:

\[
\neg P.
\]

If arguments of this sort are fundamental to dialectical
practice, the question immediately arises:

\begin{center}
What exactly does ``if \(P\), then \(Q\)'' mean?
\end{center}

\begin{whiteboard}
\textbf{ZENONIAN DIALECTIC}

\[
P\rightarrow Q
\]

\[
P\rightarrow\neg Q
\]

\[
\therefore\neg P
\]

\vspace{0.5em}

\[
\boxed{\text{WHAT DOES }P\rightarrow Q\text{ MEAN?}}
\]
\end{whiteboard}

The first extensive ancient debate over this problem appears
among Diodorus, Philo, and later dialecticians.

% ============================================================
\section{Four Ancient Theories of the Conditional}
% ============================================================

Sextus Empiricus reports four views.

The authorship of all four is not certain.

Philo and Diodorus are securely identified with the first two.

The third may well be associated with Chrysippus.

\begin{whiteboard}
\textbf{FOUR THEORIES OF ``IF \(P\), THEN \(Q\)''}

\[
\begin{array}{rl}
1.\ \textbf{Philo:}
&
\text{not }(P\text{ true and }Q\text{ false})
\\[0.5em]
2.\ \textbf{Diodorus:}
&
\text{never could / can be}\\
&
P\text{ true and }Q\text{ false}
\\[0.5em]
3.\ \textbf{Connexion:}
&
P\text{ incompatible with }\neg Q
\\[0.5em]
4.\ \textbf{Implication:}
&
Q\text{ contained potentially in }P
\end{array}
\]
\end{whiteboard}

The disagreement concerns how strong a connexion must obtain
between antecedent and consequent.

% ============================================================
\section{Philo's Truth-Functional Conditional}
% ============================================================

Philo gives the weakest of these theories.

A conditional is false only when its antecedent is true and its
consequent false.

In every other combination of truth-values, the conditional is
true.

Using modern tabular notation:

\begin{whiteboard}
\textbf{PHILO'S CONDITIONAL}

\[
\begin{array}{c|c|c}
P & Q & P\rightarrow Q\\
\hline
T & T & T\\
T & F & F\\
F & T & T\\
F & F & T
\end{array}
\]

\[
\boxed{
\text{Only false case: }
P=T,\ Q=F
}
\]
\end{whiteboard}

The truth-table device itself is modern.

The important historical fact is that the idea of
\textbf{truth-functional dependence} is already clearly present.

% ============================================================
\section{Why Philo's Theory Makes Logical Sense}
% ============================================================

At first sight Philo's theory gives strange results.

Why should:

\[
P\rightarrow Q
\]

be true merely because \(P\) is false?

The Kneales suggest that Philo may have begun from a fundamental
property required of conditional reasoning.

Whatever else ``if \(P\), then \(Q\)'' means, it must license the
following pattern:

\begin{whiteboard}
\textbf{MINIMUM REQUIREMENT FOR ``IF''}

\[
\begin{array}{c}
P\rightarrow Q\\
P\\
\hline
Q
\end{array}
\]

\[
\boxed{
\text{Philo gives the weakest truth-condition}
}
\]

\[
\boxed{
\text{that guarantees this inference.}
}
\]
\end{whiteboard}

This inference will later become the first of Chrysippus' five
basic inference schemata.

% ============================================================
\section{Diodorus' Stronger Conditional}
% ============================================================

Diodorus rejects Philo's account as too weak.

A conditional should not count as sound merely because its
antecedent happens now to be false or its consequent happens now
to be true.

Diodorus therefore adds a temporal requirement.

A sound conditional must never have been and never be capable
of beginning with a truth and ending with a falsehood.

\begin{whiteboard}
\textbf{DIODORUS}

\[
P\rightarrow Q
\]

is sound only if:

\[
\boxed{
\text{NEVER: }
P=T
\text{ and }
Q=F
}
\]

not merely:

\[
\text{not now}
\]

but:

\[
\text{across relevant times}
\]

\boardtakeaway{Diodorus strengthens Philo.}
\end{whiteboard}

This theory is closely tied to Diodorus' peculiar treatment of
modality and temporal truth.

It therefore inherits many of the difficulties of that theory.

% ============================================================
\section{The Theory of Connexion}
% ============================================================

The third reported theory says that a sound conditional must
contain a stronger connexion.

The negation of the consequent must be incompatible with the
antecedent.

Thus:

\[
P\rightarrow Q
\]

is sound when \(P\) cannot hold together with \(\neg Q\).

\begin{whiteboard}
\textbf{CONNEXION THEORY}

\[
P\rightarrow Q
\]

is sound iff:

\[
\boxed{
P\text{ is incompatible with }\neg Q
}
\]

That is:

\[
\text{not possible: }P\land\neg Q
\]
\end{whiteboard}

The surviving evidence suggests that this may have been
Chrysippus' characteristic position, but the attribution is not
certain enough to state without qualification.

% ============================================================
\section{Conditionality and Entailment}
% ============================================================

The debate exposes an important distinction.

If \(P\) logically entails \(Q\), then we may certainly form the
true conditional:

\[
P\rightarrow Q.
\]

But the converse does not automatically hold.

Not every conditional accepted as true in ordinary discourse is
therefore an expression of logical entailment.

\begin{whiteboard}
\textbf{CONDITIONALITY VS. ENTAILMENT}

\[
P\models Q
\]

supports:

\[
P\rightarrow Q
\]

Therefore:

\[
\boxed{
\text{ENTAILMENT}
\Rightarrow
\text{true conditional}
}
\]

But:

\[
\boxed{
\text{true conditional}
\not\Rightarrow
\text{logical entailment}
}
\]

\boardtakeaway{CONDITIONALITY $\neq$ ENTAILMENT}
\end{whiteboard}

Much later confusion in logic resulted from attempts to identify
all uses of ``if \ldots then \ldots'' with necessary logical
connexion.

% ============================================================
\section{The Kneales on Primitive Conditional Speech}
% ============================================================

At this point the Kneales move from historical reconstruction to
their own analysis of ordinary conditional speech.

This should be distinguished clearly from the doctrines of Philo,
Diodorus, or Chrysippus.

Consider a promise:

\begin{center}
``I will do \(Q\) if \(P\) occurs.''
\end{center}

If \(P\) never occurs, we do not ordinarily say that the speaker
has fulfilled the promise merely because the antecedent is false.

Similarly, a conditional order does not count as obeyed simply
because the relevant condition never occurs.

The Kneales suggest that one primitive function of an
``if''-clause is to select a possibility for attention.

What occurs in the consequent is then said with respect to that
possibility.

\begin{whiteboard}
\textbf{KNEALES: PRIMITIVE ``IF''}

\[
\text{If }P,\text{ then }Q
\]

Interpretation:

\[
P
\rightarrow
\text{selects a POSSIBILITY}
\]

\[
Q
\rightarrow
\text{said with respect to that possibility}
\]

If \(P\) never occurs:

\[
\boxed{\text{the conditional may have NO APPLICATION}}
\]
\end{whiteboard}

% ============================================================
\section{The Defective Truth-Function}
% ============================================================

On this analysis, a primitive conditional can be represented as
a \textbf{defective truth-function}.

When the antecedent is fulfilled, the truth or falsity of the
consequent determines whether the conditional prediction was
borne out.

When the antecedent is false, the conditional may simply have no
application.

\begin{whiteboard}
\textbf{DEFECTIVE TRUTH-FUNCTION}

\[
\begin{array}{c|c|c}
P & Q & \text{If }P\text{ then }Q\\
\hline
T & T & T\\
T & F & F\\
F & T & -\\
F & F & -
\end{array}
\]

\[
-\ =\ \text{no application}
\]

\[
\boxed{\text{Not automatically bivalent.}}
\]
\end{whiteboard}

Again, this is the Kneales' philosophical analysis, not an ancient
Stoic truth-table.

% ============================================================
\section{Necessary-Connexion Conditionals}
% ============================================================

Other conditional statements do something stronger.

The speaker may mean:

\[
P\text{ cannot occur without }Q.
\]

Such a claim is not established merely by discovering that \(P\)
and \(Q\) happen both to be true.

It asserts some modal connexion between them.

\begin{whiteboard}
\textbf{TWO IMPORTANT USES OF ``IF''}

\[
\begin{array}{rl}
\textbf{Primitive:}
&
\text{If }P,\text{ expect }Q.
\\[0.5em]
\textbf{Modal:}
&
P\text{ cannot hold without }Q.
\end{array}
\]

\[
\boxed{
\text{truth-value dependence}
\neq
\text{necessary connexion}
}
\]

\boardtakeaway{No single simple rule captures every ordinary use of ``if.''}
\end{whiteboard}

The historical importance of the Megarian debate lies partly in
the fact that ancient logicians already recognized the difficulty
of explaining conditional connexion.

% ============================================================
% PART IV
% ============================================================

\part*{IV. Stoic Meaning, Truth, and the Lekton}

% ============================================================
\section{Logic as a Part of Philosophy}
% ============================================================

The Stoics gave logic a more explicit place in their philosophical
system than any of their predecessors.

Their standard division of philosophy was:

\begin{enumerate}
    \item physics;
    \item ethics;
    \item logic.
\end{enumerate}

The Peripatetics were later represented as treating logic mainly
as an \textbf{instrument} of philosophy.

The Stoics treated it as one of philosophy's constituent parts.

\begin{whiteboard}
\textbf{STOIC PHILOSOPHY}

\[
\boxed{
\text{PHYSICS}
\qquad
\text{ETHICS}
\qquad
\text{LOGIC}
}
\]

\vspace{0.5em}

\[
\begin{array}{rcl}
\textbf{Stoics:}
&\text{logic}=&
\text{PART of philosophy}
\\[0.4em]
\textbf{Peripatetics:}
&\text{logic}=&
\text{INSTRUMENT of philosophy}
\end{array}
\]
\end{whiteboard}

Within logic, the Stoic term \textbf{dialectic} corresponds most
closely to what we now call logic, though Stoic dialectic also
includes material we would place under grammar, semantics, and
epistemology.

% ============================================================
\section{Things Signifying and Things Signified}
% ============================================================

The Stoics divided dialectic into investigation of:

\begin{itemize}
    \item things that signify;
    \item things that are signified.
\end{itemize}

They developed a remarkably sophisticated theory of linguistic
meaning.

The central distinction involves three items:

\begin{enumerate}
    \item the signifying expression;
    \item what is signified or meant;
    \item the external object.
\end{enumerate}

What is signified or meant is called a \textbf{lekton}.

\begin{whiteboard}
\textbf{STOIC SEMANTIC TRIAD}

\[
\begin{array}{c}
\textbf{SIGNIFIER}\\
\text{spoken expression}\\
\downarrow\\
\textbf{LEKTON}\\
\text{what is meant / signified}\\
\downarrow\\
\textbf{OBJECT}\\
\text{thing in the world}
\end{array}
\]
\end{whiteboard}

In Stoic metaphysics:

\begin{itemize}
    \item speech is corporeal;
    \item the external object is corporeal;
    \item the lekton is incorporeal.
\end{itemize}

\begin{whiteboard}
\[
\begin{array}{rcl}
\text{speech}
&=&
\text{corporeal}
\\[0.3em]
\text{object}
&=&
\text{corporeal}
\\[0.3em]
\text{lekton}
&=&
\boxed{\text{INCORPOREAL}}
\end{array}
\]
\end{whiteboard}

This is striking because the Stoics were otherwise strongly
materialist philosophers.

% ============================================================
\section{Why the Lekton Is Logically Important}
% ============================================================

The theory responds to a problem that has appeared repeatedly
throughout our course.

What exactly is true or false?

A physical sound cannot by itself be the primary bearer of truth.

Someone who does not understand Greek can hear a Greek
sentence without understanding what it signifies.

The Stoics therefore distinguish the expression from what the
expression means.

\begin{whiteboard}
\textbf{OLD PROBLEM}

What is TRUE or FALSE?

Not merely:

\[
\text{sound}
\]

Not merely:

\[
\text{physical inscription}
\]

Stoic distinction:

\[
\boxed{
\text{speech}
\rightarrow
\text{expresses}
\rightarrow
\text{LEKTON}
}
\]
\end{whiteboard}

But we must not immediately identify a lekton with a modern
proposition.

The Stoic theory is more complicated.

% ============================================================
\section{Complete and Defective Lekta}
% ============================================================

The Stoics divide lekta into complete and defective.

A defective lekton does not by itself constitute a finished
content.

For example:

\begin{center}
``writes''
\end{center}

leaves us asking:

\begin{center}
Who writes?
\end{center}

By contrast:

\begin{center}
``Socrates writes''
\end{center}

expresses a complete lekton.

\begin{whiteboard}
\textbf{LEKTA}

\[
\begin{array}{rl}
\textbf{DEFECTIVE:}
&
\text{``\ldots writes''}
\\
&
\Rightarrow \text{Who?}
\\[0.8em]
\textbf{COMPLETE:}
&
\text{``Socrates writes.''}
\end{array}
\]

\[
\text{predicate alone}
\rightarrow
\text{incomplete}
\]

\[
\text{finished content}
\rightarrow
\text{complete}
\]
\end{whiteboard}

Defective lekta include such things as predicates and subjects
considered apart from complete assertion.

% ============================================================
\section{Axiomata}
% ============================================================

Among complete lekta, the most important logical class is the
\textbf{axioma}.

The word must not be confused with the modern mathematical word
``axiom.''

A Stoic axioma is roughly a complete declarative content capable
of truth or falsity.

Other complete lekta include:

\begin{itemize}
    \item questions;
    \item commands;
    \item prayers;
    \item oaths;
    \item suppositions;
    \item addresses.
\end{itemize}

These need not themselves be true or false.

\begin{whiteboard}
\textbf{COMPLETE LEKTA}

\[
\begin{array}{rcl}
\textbf{AXIOMA}
&\rightarrow&
\text{TRUE or FALSE}
\\[0.4em]
\text{question}
&\rightarrow&
\text{not true / false}
\\[0.4em]
\text{command}
&\rightarrow&
\text{not true / false}
\\[0.4em]
\text{prayer}
&\rightarrow&
\text{not true / false}
\end{array}
\]

\boardtakeaway{AXIOMA = complete declaratory lekton capable of truth or falsity.}
\end{whiteboard}

This is a significant advance in the analysis of meaning.

Not everything that has complete meaning is therefore a bearer of
truth-value.

% ============================================================
\section{Where Truth and Falsity Belong}
% ============================================================

The Stoics explicitly distinguish several candidates for truth and
falsity.

According to the doctrine reported by Sextus, the fundamental
bearer of truth is neither:

\begin{itemize}
    \item the spoken expression as such;
    \item nor merely a physical or psychological movement of
    thought.
\end{itemize}

Truth and falsity fundamentally concern what is signified.

For logical purposes, the axioma becomes the crucial bearer of
truth-value.

\begin{whiteboard}
\textbf{WHAT IS TRUE OR FALSE?}

Not fundamentally:

\[
\text{sound / sentence-token}
\]

Not fundamentally:

\[
\text{mental movement}
\]

Stoic answer:

\[
\boxed{
\text{AXIOMA / SIGNIFIED CONTENT}
}
\]

\[
\text{truth is fundamentally attached to what is meant}
\]
\end{whiteboard}

The Stoics also used ``true'' and ``false'' in derivative ways for
presentations and arguments.

A true argument, for example, was valid and had true premisses.

% ============================================================
\section{The True and the Truth}
% ============================================================

The Stoics even distinguish:

\begin{center}
\textbf{the true}
\end{center}

from:

\begin{center}
\textbf{truth}.
\end{center}

The true is an incorporeal axioma.

Truth, in the special Stoic sense, is a body of knowledge
possessed by a knower.

This distinction belongs partly to Stoic metaphysics and
epistemology rather than formal logic, so we need only record its
structure.

\begin{whiteboard}
\textbf{STOIC DISTINCTION}

\[
\begin{array}{rcl}
\textbf{the true}
&=&
\text{a true axioma}
\\
&&
\text{incorporeal}
\\[0.7em]
\textbf{truth}
&=&
\text{a body of knowledge}
\\
&&
\text{corporeal in Stoic theory}
\end{array}
\]
\end{whiteboard}

Our logical concern remains chiefly with true and false axiomata.

% ============================================================
\section{Negation and the Discovery of Scope}
% ============================================================

The Stoics give careful attention to negation.

To form the contradictory of an axioma, the negative must govern
the \textbf{whole} axioma.

Compare:

\[
\neg(P\land Q)
\]

with:

\[
P\land\neg Q.
\]

These are not the same proposition.

In the first, negation applies to the entire conjunction.

In the second, it applies only to \(Q\).

\begin{whiteboard}
\textbf{SCOPE OF NEGATION}

\[
\boxed{\neg(P\land Q)}
\]

means:

\[
\text{NOT: both }P\text{ and }Q
\]

But:

\[
P\land\neg Q
\]

means:

\[
P\text{ and not-}Q
\]

Therefore:

\[
\boxed{
\neg(P\land Q)
\neq
P\land\neg Q
}
\]

\boardtakeaway{Negation of the whole $\neq$ negation of one part.}
\end{whiteboard}

The Kneales identify this discussion as the first appearance in
the history of logic of something recognizably like the notion of
the \textbf{scope of an operator}.

% ============================================================
\section{Double Negation}
% ============================================================

The Stoics also recognize a double negative and hold it equivalent
to the corresponding affirmative.

In modern notation:

\[
\neg\neg P
\]

is equivalent to:

\[
P.
\]

\begin{whiteboard}
\textbf{DOUBLE NEGATION}

\[
\boxed{
\neg\neg P
\Longleftrightarrow
P
}
\]

Stoic terminology includes the idea of a
``super-negative.''
\end{whiteboard}

This principle later becomes important in the reconstruction of
Stoic derivations.

% ============================================================
\section{Universal Statements as Conditionals}
% ============================================================

Another striking difference from Aristotle appears in the Stoic
treatment of general propositions.

A statement such as:

\begin{center}
Man is a rational mortal animal
\end{center}

can be represented approximately as:

\begin{center}
If anything is a man, then it is a rational mortal animal.
\end{center}

Thus universality is expressed by conditional structure.

\begin{whiteboard}
\textbf{ARISTOTELIAN STYLE}

\[
\text{Every man is mortal.}
\]

\textbf{STOIC-STYLE ANALYSIS}

\[
\boxed{
\text{If anything is a man, then it is mortal.}
}
\]

\[
\text{universal predication}
\rightarrow
\text{conditional structure}
\]
\end{whiteboard}

This helps explain why conditionals become so central in Stoic
logic.

% ============================================================
\section{Simple and Complex Axiomata}
% ============================================================

The Stoics distinguish simple from non-simple or complex
axiomata.

Complex axiomata are constructed by means of connectives.

The ancient classification includes several forms.

For our logical purposes the principal structures can be displayed
using modern notation.

\begin{whiteboard}
\textbf{COMPLEX AXIOMATA}

\[
\begin{array}{rcl}
\neg P
&=&
\text{negation}
\\[0.4em]
P\land Q
&=&
\text{conjunction}
\\[0.4em]
P\lor Q
&=&
\text{disjunction}
\\[0.4em]
P\rightarrow Q
&=&
\text{conditional}
\\[0.4em]
\neg(P\land Q)
&=&
\text{``not both''}
\end{array}
\]

\boardtakeaway{Modern symbols are ours, not Stoic notation.}
\end{whiteboard}

The Stoic classification also includes causal and comparative
complexes, but these need not occupy our main formal discussion.

% ============================================================
\section{Conjunction}
% ============================================================

The Stoics recognized the familiar truth condition for
conjunction.

A conjunction is true when all of its components are true and
false when at least one component is false.

\begin{whiteboard}
\textbf{CONJUNCTION}

\[
P\land Q
\]

is TRUE iff:

\[
P=T
\quad\text{and}\quad
Q=T
\]

Otherwise:

\[
P\land Q=F
\]
\end{whiteboard}

This is a clear example of truth-functional treatment.

% ============================================================
\section{Stoic Disjunction}
% ============================================================

The early Stoic disjunction is normally strong or exclusive.

Thus:

\[
P\xor Q
\]

means:

\begin{center}
one or the other, but not both.
\end{center}

Later Stoics also recognized a weaker form under which the
disjunction is true if at least one disjunct is true.

\begin{whiteboard}
\textbf{STOIC ``OR''}

Early strong form:

\[
P\xor Q
\]

means:

\[
\boxed{
\text{one OR the other, BUT NOT BOTH}
}
\]

Later tradition also recognizes:

\[
P\lor Q
\]

where both may be true.
\end{whiteboard}

This distinction is important because Chrysippus' fourth and fifth
indemonstrables depend upon the strong disjunction.

% ============================================================
\section{Axioma and the Modern Proposition}
% ============================================================

It is tempting to identify:

\[
\text{axioma}
=
\text{proposition}.
\]

The comparison is useful but not exact.

Stoic axiomata can possess features that modern philosophers
often exclude from abstract propositions.

They can involve:

\begin{itemize}
    \item mood;
    \item tense;
    \item context-sensitive reference;
    \item changing truth-values.
\end{itemize}

Some Stoic discussions even allow an axioma to be
``destroyed'' when the mode of reference upon which it depends
can no longer operate.

\begin{whiteboard}
\textbf{AXIOMA $\approx$ PROPOSITION}

But not exactly.

Stoic axioma may have:

\[
\begin{array}{c}
\text{mood}\\
\text{tense}\\
\text{context}\\
\text{changing truth-value}
\end{array}
\]

and may even:

\[
\text{cease to be expressible}
\]

\[
\downarrow
\]

\[
\text{``destroyed''}
\]

\boardtakeaway{Do not turn the Stoics into modern logicians too quickly.}
\end{whiteboard}

The Stoics have clearly distinguished linguistic expression from
what it means, but their lekta preserve more features of
utterance-context than many later theories of propositions do.

% ============================================================
% PART V
% ============================================================

\part*{V. The Stoic System of Inference Schemata}

% ============================================================
\section{The Fundamental Contrast with Aristotle}
% ============================================================

We have now reached the formal center of Stoic logic.

The key contrast with Aristotle concerns the kinds of variables
used.

Aristotle uses letters as places for \textbf{terms}.

The Stoics instead use ordinal expressions such as:

\begin{center}
the first, the second, the third
\end{center}

as places for \textbf{whole propositions}.

This difference is fundamental.

\begin{whiteboard}
\textbf{ARISTOTLE}

\[
\begin{array}{l}
\text{Every }M\text{ is }L\\
\text{Every }S\text{ is }M\\
\hline
\therefore \text{Every }S\text{ is }L
\end{array}
\]

\[
M,S,L
\rightarrow
\boxed{\text{TERM VARIABLES}}
\]
\end{whiteboard}

\begin{whiteboard}
\textbf{STOICS}

\[
\begin{array}{l}
\text{If the FIRST, then the SECOND}\\
\text{The FIRST}\\
\hline
\therefore \text{The SECOND}
\end{array}
\]

\[
P,Q
\rightarrow
\boxed{\text{PROPOSITIONAL VARIABLES}}
\]
\end{whiteboard}

This is the key to understanding why Stoic logic anticipates what
later came to be called propositional logic.

% ============================================================
\section{Inference Schemata}
% ============================================================

Where Aristotle often formulates a logical principle as a
conditional statement about terms, the Stoics regularly present
\textbf{skeleton arguments}.

Using modern letters merely for convenience:

\begin{whiteboard}
\textbf{INFERENCE SCHEMA}

\[
\begin{array}{c}
P\rightarrow Q\\
P\\
\hline
Q
\end{array}
\]

\[
\boxed{
P,Q=\text{whole propositions}
}
\]

\boardtakeaway{The form is a pattern of ARGUMENT, not merely a relation among terms.}
\end{whiteboard}

The Stoics themselves used ordinal expressions rather than our
letters \(P,Q,R\).

% ============================================================
\section{The Five Indemonstrables}
% ============================================================

Chrysippus recognizes five fundamental inference schemata.

Ancient authors call them the
\textbf{indemonstrable moods}.

They function as basic inference patterns from which more
complicated arguments can be analyzed.

Because the board is small, they should be written in two
packets.

% ============================================================
\subsection{First and Second Indemonstrables}
% ============================================================

\begin{whiteboard}
\textbf{FIVE INDEMONSTRABLES --- I}

\textbf{1.}

\[
\begin{array}{c}
P\rightarrow Q\\
P\\
\hline
Q
\end{array}
\]

\textbf{2.}

\[
\begin{array}{c}
P\rightarrow Q\\
\neg Q\\
\hline
\neg P
\end{array}
\]

Modern later names:

\[
1=\text{modus ponens}
\]

\[
2=\text{modus tollens}
\]

\textit{Copy, then erase.}
\end{whiteboard}

The Latin names are later.

They are useful labels for us, but they should not be projected
back onto Chrysippus.

% ============================================================
\subsection{Third, Fourth, and Fifth Indemonstrables}
% ============================================================

The third uses a ``not both'' proposition.

The fourth and fifth use the Stoics' strong or exclusive
disjunction.

\begin{whiteboard}
\textbf{FIVE INDEMONSTRABLES --- II}

\textbf{3.}

\[
\begin{array}{c}
\neg(P\land Q)\\
P\\
\hline
\neg Q
\end{array}
\]

\textbf{4.}

\[
\begin{array}{c}
P\xor Q\\
P\\
\hline
\neg Q
\end{array}
\]

\textbf{5.}

\[
\begin{array}{c}
P\xor Q\\
\neg Q\\
\hline
P
\end{array}
\]

\[
\xor
=
\text{exclusive OR}
\]
\end{whiteboard}

Together these five patterns form the basic machinery of
Chrysippean inference.

% ============================================================
\section{Why They Are Called Indemonstrable}
% ============================================================

``Indemonstrable'' does not mean unintelligible or irrational.

The five forms function as basic inference patterns rather than
patterns that must themselves be derived from still more
fundamental Stoic moods.

More complicated valid forms can then be analyzed in terms of
them.

\begin{whiteboard}
\textbf{INDEMONSTRABLES}

\[
\boxed{\text{BASIC inference patterns}}
\]

\[
\begin{array}{c}
\text{five basic moods}\\
\downarrow\\
\text{formal derivation}\\
\downarrow\\
\text{more complex moods}
\end{array}
\]

\boardtakeaway{Logic becomes a deductive system of inference rules.}
\end{whiteboard}

This is a crucial difference between merely collecting examples
of good arguments and constructing a logical system.

% ============================================================
\section{Leading Premise and Additional Assumption}
% ============================================================

For Chrysippus, the simplest argument normally contains two
premisses.

One is complex and determines the characteristic structure of the
argument.

This is called the \textbf{leading premiss}.

The second is the \textbf{additional assumption}.

\begin{whiteboard}
\textbf{STOIC BASIC ARGUMENT}

\[
\boxed{\text{Leading premiss}}
\]

\[
\text{complex proposition}
\]

\[
+
\]

\[
\boxed{\text{Additional assumption}}
\]

\[
\downarrow
\]

\[
\text{CONCLUSION}
\]
\end{whiteboard}

Chrysippus regarded the two-premiss form as basic.

Antipater later recognized one-premiss inference, and this was
remembered as an innovation.

\begin{whiteboard}
\[
\begin{array}{rcl}
\text{Chrysippus:}
&
\text{simplest inference}
&=
2\text{ premisses}
\\[0.4em]
\text{Antipater:}
&
\text{one-premiss inference}
&=
\text{allowed}
\end{array}
\]
\end{whiteboard}

% ============================================================
\section{Derived Inference}
% ============================================================

The most important feature of the Stoic system is that complicated
arguments can be reduced to sequences of simpler ones.

Consider:

\[
(P\land Q)\rightarrow R,
\]

\[
\neg R,
\]

\[
P.
\]

We wish to derive:

\[
\neg Q.
\]

First use the second indemonstrable on:

\[
(P\land Q)\rightarrow R
\]

and:

\[
\neg R.
\]

This yields:

\[
\neg(P\land Q).
\]

Then apply the third indemonstrable to:

\[
\neg(P\land Q)
\]

and:

\[
P.
\]

The result is:

\[
\neg Q.
\]

Because the board is small, do the derivation in packets.

\begin{whiteboard}
\textbf{DERIVED ARGUMENT --- SETUP}

\[
\begin{array}{c}
(P\land Q)\rightarrow R\\
\neg R\\
P\\
\hline
\therefore \neg Q
\end{array}
\]

\textit{Copy, then erase.}
\end{whiteboard}

\begin{whiteboard}
\textbf{DERIVED ARGUMENT --- STEP 1}

\[
\begin{array}{c}
(P\land Q)\rightarrow R\\
\neg R\\
\hline
\neg(P\land Q)
\end{array}
\]

Using Indemonstrable 2.
\end{whiteboard}

\begin{whiteboard}
\textbf{DERIVED ARGUMENT --- STEP 2}

\[
\begin{array}{c}
\neg(P\land Q)\\
P\\
\hline
\neg Q
\end{array}
\]

Using Indemonstrable 3.

\[
\boxed{
\text{Complex inference}
=
\text{sequence of basic moods}
}
\]
\end{whiteboard}

This is one of the clearest surviving examples of the systematic
character of Stoic logic.

% ============================================================
\section{Intermediate Conclusions}
% ============================================================

The preceding demonstration depends on an important proof
principle.

What has been derived from some premisses may itself be used as
a premiss in a further inference.

Schematically:

\[
P_1,P_2\vdash Q
\]

and:

\[
Q,P_3\vdash R.
\]

Then:

\[
P_1,P_2,P_3\vdash R.
\]

\begin{whiteboard}
\textbf{INTERMEDIATE CONCLUSIONS}

\[
\begin{array}{c}
P_1,P_2\vdash Q\\
Q,P_3\vdash R\\
\hline
P_1,P_2,P_3\vdash R
\end{array}
\]

\boardtakeaway{An intermediate conclusion may become a new premiss.}
\end{whiteboard}

This principle is obvious to a modern student familiar with
formal proof, but its explicit use is historically significant.

% ============================================================
\section{The Themata}
% ============================================================

Ancient reports also tell us that the Stoics used four
\textbf{themata} in the analysis of complicated arguments.

Unfortunately, only fragments of this doctrine survive.

The first and third are known with reasonable clarity.

The second and fourth are lost and can only be reconstructed
conjecturally.

The Kneales therefore treat this part of Stoic logic cautiously.

\begin{whiteboard}
\textbf{THEMATA}

\[
\boxed{
\text{general rules for ANALYZING complex arguments}
}
\]

Four were recognized.

But:

\[
\text{1st + 3rd survive clearly}
\]

\[
\text{2nd + 4th are lost}
\]
\end{whiteboard}

% ============================================================
\section{The First Thema}
% ============================================================

The first thema can be represented approximately as follows.

If two premisses entail a conclusion, then either premiss together
with the negation of the conclusion entails the negation of the
other premiss.

Schematically:

\[
P,Q\vdash R
\]

gives:

\[
P,\neg R\vdash\neg Q.
\]

\begin{whiteboard}
\textbf{FIRST THEMA}

If:

\[
P,Q\vdash R
\]

then:

\[
\boxed{
P,\neg R\vdash\neg Q
}
\]

and correspondingly:

\[
Q,\neg R\vdash\neg P
\]
\end{whiteboard}

This is closely related to the familiar technique of indirect
argument.

% ============================================================
\section{The Third Thema}
% ============================================================

The third thema concerns the substitution of a derivation into a
larger argument.

Suppose some argument requires \(Q\) as a premiss.

If \(Q\) has itself been established from \(P_1\) and \(P_2\),
those supporting premisses can replace \(Q\) in the larger
derivation.

\begin{whiteboard}
\textbf{THIRD THEMA}

Suppose:

\[
P_1,P_2\vdash Q
\]

and:

\[
Q,R\vdash S.
\]

Then:

\[
\boxed{
P_1,P_2,R\vdash S
}
\]

\[
\text{derived premiss}
\rightarrow
\text{replace by its supporting premisses}
\]
\end{whiteboard}

The principle helps explain how the Stoics could reduce
complicated arguments to combinations of their basic forms.

% ============================================================
\section{Conditionalization}
% ============================================================

The Kneales' reconstruction of Stoic procedure suggests another
important principle.

If:

\[
P,Q\vdash R,
\]

the corresponding necessary connexion can be represented by the
conditional:

\[
(P\land Q)\rightarrow R.
\]

This allows a valid inference relation to be introduced as a
conditional premiss inside another argument.

\begin{whiteboard}
\textbf{CONDITIONALIZATION}

\[
P,Q\vdash R
\]

becomes:

\[
\boxed{
(P\land Q)\rightarrow R
}
\]

\[
\text{inference pattern}
\rightarrow
\text{conditional premiss}
\]
\end{whiteboard}

This is an important modern reconstruction of how the Stoic
system may have operated.

It should not be presented as though this exact symbolic rule
survived in Chrysippus' own writings.

% ============================================================
\section{Reductio Returns}
% ============================================================

One of the most important Stoic derived patterns brings us back
to the very origins of the Megarian tradition.

Suppose:

\[
P\rightarrow Q
\]

and:

\[
P\rightarrow\neg Q.
\]

Then:

\[
\neg P.
\]

\begin{whiteboard}
\textbf{STOIC DERIVED REDUCTIO}

\[
\begin{array}{c}
P\rightarrow Q\\
P\rightarrow\neg Q\\
\hline
\neg P
\end{array}
\]

Historical line:

\[
\text{Zeno}
\rightarrow
\text{Megarians}
\rightarrow
\text{Chrysippus}
\]

\boardtakeaway{Dialectical practice has become formal theory.}
\end{whiteboard}

This is historically beautiful.

The style of argument that had helped stimulate Megarian logic
has now been incorporated into a systematic formal calculus.

% ============================================================
\section{Consequentia Mirabilis}
% ============================================================

Another striking form associated with Stoic logical development
is:

\[
\neg P\rightarrow P.
\]

From this, under the relevant principles of the system, one may
derive:

\[
P.
\]

The pattern later came to be known as
\textit{consequentia mirabilis}.

\begin{whiteboard}
\textbf{CONSEQUENTIA MIRABILIS}

\[
\begin{array}{c}
\neg P\rightarrow P\\
\hline
P
\end{array}
\]

Interpretive idea:

\[
\boxed{
\text{If denying }P\text{ forces }P,\text{ then }P.
}
\]
\end{whiteboard}

The pattern became important in later anti-sceptical and logical
arguments.

But we should not assume that every philosophical application of
the form establishes as much as its user claims.

% ============================================================
\section{Rigour and Formalism}
% ============================================================

Ancient Peripatetics sometimes criticized the Stoics for excessive
attention to verbal or formal details.

Chrysippus apparently proved even results that other philosophers
were content simply to assume.

To modern students trained in mathematical reasoning, this
``excessive'' concern for form may look much more like a virtue
than a defect.

\begin{whiteboard}
\textbf{ANCIENT CRITICISM}

Stoics:

\[
\text{too formal}
\]

\[
\text{too concerned with exact expression}
\]

Modern perspective:

\[
\boxed{
\text{formal precision is a major logical achievement}
}
\]
\end{whiteboard}

Unfortunately, the hostile ancient reputation of Stoic formalism
may have contributed to the disappearance of Chrysippus'
logical works.

% ============================================================
\section{Did the Stoics Claim Completeness?}
% ============================================================

Ancient reports suggest that Chrysippus believed complicated valid
arguments could ultimately be analyzed through the five
indemonstrables.

This resembles a claim of completeness.

But we must be careful.

The Stoics did not possess our modern technical definition of a
complete formal calculus, nor do we possess an ancient proof of
completeness in the modern sense.

\begin{whiteboard}
\textbf{STOIC SYSTEMATIC CLAIM}

\[
\text{valid arguments}
\]

\[
\downarrow
\]

\[
\text{analysis}
\]

\[
\downarrow
\]

\[
\boxed{\text{FIVE INDEMONSTRABLES}}
\]

But:

\[
\boxed{
\text{ancient ``complete''}
\neq
\text{modern completeness theorem}
}
\]
\end{whiteboard}

Modern reconstructions can prove strong results for reconstructed
Stoic systems, but those achievements must not simply be
attributed to Chrysippus without evidence.

% ============================================================
\section{Aristotle and Chrysippus Compared}
% ============================================================

We can now make the central comparison of the chapter.

Aristotle's variables range over \textbf{general terms}.

Stoic variables range over \textbf{propositions}.

Aristotle therefore analyzes logical relations involving the
internal structure of propositions.

The Stoics analyze patterns of inference determined by the ways
whole propositions are joined by connectives.

\begin{whiteboard}
\textbf{ARISTOTLE}

\[
\begin{array}{l}
\text{Every }M\text{ is }L\\
\text{Every }S\text{ is }M\\
\hline
\text{Every }S\text{ is }L
\end{array}
\]

\[
\boxed{\text{variables } \rightarrow \text{ TERMS}}
\]
\end{whiteboard}

\begin{whiteboard}
\textbf{CHRYSIPPUS}

\[
\begin{array}{c}
P\rightarrow Q\\
P\\
\hline
Q
\end{array}
\]

\[
\boxed{\text{variables } \rightarrow \text{ PROPOSITIONS}}
\]
\end{whiteboard}

The distinction can be stated in one line.

\begin{whiteboard}
\[
\begin{array}{rcl}
\textbf{Aristotle}
&\rightarrow&
\text{relations WITHIN propositions}
\\[0.6em]
\textbf{Stoics}
&\rightarrow&
\text{relations BETWEEN propositions}
\end{array}
\]
\end{whiteboard}

This is why the two systems are complementary rather than
genuine competitors.

% ============================================================
\section{Primary Logic}
% ============================================================

The Kneales propose the term \textbf{primary logic} for the
fundamental theory involving such notions as:

\[
\neg,
\qquad
\land,
\qquad
\lor,
\qquad
\rightarrow.
\]

This logic deals with the combinations and inferential relations
of propositions without yet introducing the generality expressed
by ``every'' and ``some.''

\begin{whiteboard}
\textbf{PRIMARY LOGIC}

\[
\boxed{
\neg
\qquad
\land
\qquad
\lor
\qquad
\rightarrow
}
\]

Concern:

\[
\text{relations among propositions}
\]
\end{whiteboard}

Stoic dialectic is the first substantial ancient version of this
kind of logic.

% ============================================================
\section{General Logic}
% ============================================================

The Kneales reserve the expression \textbf{general logic} for the
larger theory in which primary logic is combined with notions of
generality such as:

\[
\text{every}
\]

and:

\[
\text{some}.
\]

Thus general logic includes quantificational structure.

\begin{whiteboard}
\textbf{GENERAL LOGIC}

\[
\boxed{
\text{PRIMARY LOGIC}
+
\text{EVERY / SOME}
}
\]

That is:

\[
\text{connectives}
+
\text{quantification}
\]
\end{whiteboard}

Primary logic is more fundamental not because it contains general
logic, but because general logic presupposes it.

\begin{whiteboard}
\[
\boxed{
\text{General logic presupposes primary logic.}
}
\]

But:

\[
\text{primary logic can be developed independently.}
\]
\end{whiteboard}

% ============================================================
\section{Where Aristotle and Chrysippus Fit}
% ============================================================

With this terminology, the Kneales give a striking reinterpretation
of ancient logic.

Aristotle's syllogistic is a \textbf{fragment of general logic}.

It deals explicitly with generality and relations among general
terms.

But it tacitly relies upon principles of primary logic that
Aristotle does not formulate as a separate system.

Chrysippus, by contrast, produces the first major ancient
\textbf{primary logic}.

\begin{whiteboard}
\textbf{TWO ACHIEVEMENTS}

\[
\begin{array}{rcl}
\textbf{Chrysippus}
&\rightarrow&
\boxed{\text{first primary logic}}
\\[0.8em]
\textbf{Aristotle}
&\rightarrow&
\boxed{\text{fragment of general logic}}
\end{array}
\]

Aristotle:

\[
\text{quantification without a separate}
\]

\[
\text{formal calculus of propositional connexion}
\]

Chrysippus:

\[
\text{propositional connexion without}
\]

\[
\text{Aristotle's quantificational term theory}
\]
\end{whiteboard}

The two systems solve different parts of the larger problem.

% ============================================================
\section{The Correct Historical Picture}
% ============================================================

Ancient partisans tended to ask:

\begin{center}
Aristotle or Chrysippus?
\end{center}

The more fruitful question is:

\begin{center}
How do the two systems fit together?
\end{center}

\begin{whiteboard}
\textbf{NOT:}

\[
\text{ARISTOTLE}
\quad
\mathbf{VS.}
\quad
\text{CHRYSIPPUS}
\]

\vspace{0.6em}

\textbf{BUT:}

\[
\boxed{
\text{ARISTOTLE}
+
\text{CHRYSIPPUS}
}
\]

\[
\boxed{
\text{TERM / GENERAL LOGIC}
+
\text{PROPOSITIONAL / PRIMARY LOGIC}
}
\]
\end{whiteboard}

The intellectual loss of antiquity was not that one school chose
the wrong logic.

It was that the complementary character of the two achievements
was obscured by rivalry.

% ============================================================
\section{Chapter Summary}
% ============================================================

The Megarian--Stoic tradition begins from a different source than
Aristotle's logic.

Instead of taking geometrical demonstration as its principal
model, it grows from dialectic, reductio, paradox, and argument
with complex propositions.

Eubulides' paradoxes expose problems of truth, identity,
vagueness, and presupposition.

Diodorus and Philo develop rival accounts of possibility.

Diodorus' Master Argument attempts to tie possibility to what is
or will be true, but the Kneales argue that the reasoning depends
upon confusion between sentences, tense, and the contents those
sentences express.

The Megarian debate over conditionals then asks what makes an
``if \ldots then \ldots'' proposition sound.

Philo develops a truth-functional account.

Diodorus proposes a stronger temporally constrained account.

The theory of connexion requires incompatibility between the
antecedent and the negation of the consequent.

The Stoics then develop a sophisticated semantic theory
distinguishing expressions, objects, and incorporeal lekta.

Among complete lekta, axiomata are the principal bearers of
truth and falsity.

Stoic discussion of negation introduces an explicit concern with
the scope of an operator.

Finally, Chrysippus systematizes inference among whole
propositions through five indemonstrable moods and procedures
for deriving more complex arguments.

\begin{whiteboard}
\textbf{CHAPTER III SUMMARY}

\[
\begin{array}{c}
\text{Megarian paradox}\\
\downarrow\\
\text{logical problems}
\\[0.8em]
\text{modality}\\
\downarrow\\
\text{Diodorus / Philo}
\\[0.8em]
\text{conditionals}\\
\downarrow\\
\text{truth-function / connexion}
\\[0.8em]
\text{lekta + axiomata}\\
\downarrow\\
\text{meaning / truth / scope}
\\[0.8em]
\text{five indemonstrables}\\
\downarrow\\
\text{systematic propositional inference}
\end{array}
\]
\end{whiteboard}

% ============================================================
\section{Final Historical Takeaway}
% ============================================================

The great achievement of Aristotle was the systematic study of
inferences involving general terms.

The great achievement of the Stoics was the systematic study of
inferences whose validity depends upon the logical relations
among whole propositions.

Neither achievement eliminates the other.

A mature logic requires both.

\begin{whiteboard}
\textbf{FINAL TAKEAWAY}

\[
\begin{array}{ccc}
\textbf{ARISTOTLE}
&&
\textbf{CHRYSIPPUS}
\\
\text{terms}
&&
\text{propositions}
\\
\text{syllogistic}
&&
\text{connectives}
\\
\text{general logic}
&&
\text{primary logic}
\end{array}
\]

\vspace{0.8em}

\[
\boxed{
\begin{array}{c}
\text{Ancient logic discovered two major pieces}\\
\text{of one larger logical system.}
\end{array}
}
\]

\vspace{0.8em}

\[
\text{Next: Roman and Medieval Logic}
\]
\end{whiteboard}

\end{document}